\section{Dati}\label{sec:data}

\subsection{ISTAT: Flussi di Emigrazione dal Lato Italiano}
\label{subsec:data_istat}
Il lato italiano del confronto è tratto dal dataset \textit{Cancellazioni anagrafiche per trasferimento di residenza all'estero}, pubblicato annualmente dall'ISTAT~\cite{istat_iscrizioni}.
Ogni record rappresenta un individuo che ha formalmente rimosso la propria residenza da un comune italiano e dichiarato un paese straniero come nuova dimora.
Questo meccanismo di misurazione è strutturalmente legato a una decisione umana: una persona deve avviare attivamente la procedura di cancellazione, sia di persona presso il proprio comune di origine sia, dal 2012, attraverso il portale online dell'AIRE.
Di conseguenza, la serie ISTAT non misura l'emigrazione in sé, ma la frazione \emph{amministrativamente visibile} dell'emigrazione: quegli individui che hanno scelto di ufficializzare la propria partenza.
Per questo studio, abbiamo filtrato il dataset per mantenere solo i record relativi alla cittadinanza italiana, entrambi i sessi, tutte le classi di età, coprendo gli anni dal 2017 al 2021 e gli otto paesi di destinazione per i quali è disponibile una corrispondente osservazione Eurostat.
\subsection{Eurostat: Flussi di Immigrazione dal Lato della Destinazione}
\label{subsec:data_eurostat}

La prospettiva del paese di destinazione è tratta dal dataset Eurostat \texttt{MIGR\_IMM1CTZ} (\textit{Immigration by age group, sex and citizenship})~\cite{eurostat_migr_imm1ctz}.
Questo dataset aggrega i record di immigrazione inviati da ciascuno stato membro a Eurostat ai sensi del Regolamento (CE) n. 862/2007~\cite{ec_regulation_862_2007}, che standardizza la definizione di migrante a lungo termine come una persona che stabilisce la residenza per almeno dodici mesi.
Per ogni coppia paese-anno, abbiamo estratto il conteggio totale dei cittadini italiani registrati come nuovi immigrati.
Il dataset copre otto paesi di destinazione: Austria, Danimarca, Finlandia, Francia, Germania, Norvegia, Spagna e Svezia, per il periodo dal 2017 al 2021.

\subsection{Costruzione e Allineamento del Dataset}
\label{subsec:data_alignment}

Le due serie sono allineate sulla chiave (paese, anno), producendo un panel di 40 osservazioni su 8 paesi e 5 anni (2017--2021).
Per ogni cella, calcoliamo il gap grezzo come la differenza tra il conteggio Eurostat e il conteggio ISTAT.
Due asimmetrie definitorie meritano un riconoscimento esplicito.
In primo luogo, l'ISTAT registra l'anno di cancellazione, mentre Eurostat registra l'anno di arrivo; per gli individui che si cancellano con ritardo, le due cifre possono cadere in anni solari diversi.
In secondo luogo, la disaggregazione per età non è disponibile in modo coerente tra i vari paesi in \texttt{MIGR\_IMM1CTZ}; di conseguenza, l'analisi copre l'intera popolazione emigrante.
\paragraph{Garantire la comparabilità della popolazione.}
Il registro delle cancellazioni ISTAT registra \emph{tutti} gli individui che si cancellano da un comune italiano, indipendentemente dalla cittadinanza: questo include cittadini stranieri precedentemente residenti in Italia che si trasferiscono all'estero.
Abbiamo affrontato questo problema limitando l'estrazione ISTAT ai soli record con cittadinanza italiana, mantenendo solo quegli individui per i quali la decisione di iscrizione all'AIRE è effettivamente pertinente.
Sul lato Eurostat, il dataset \texttt{MIGR\_IMM1CTZ} è nativamente filtrato per cittadinanza dell'immigrante.
Dopo questa fase, entrambe le serie misurano la stessa popolazione sottostante: cittadini italiani che attraversano un confine internazionale e stabiliscono la propria residenza all'estero.
